\section{Проверка гипотезы о параметре нормального распределения \(X_3\)}

\subsection{Постановка задачи}

По столбцу \(X_3\) проверяется гипотеза о математическом ожидании нормального распределения (дисперсия неизвестна).

\subsubsection{Гипотезы}
\begin{itemize}
    \item \(H_0: \mu = \mu_0 = 75.24\) — математическое ожидание \(X_3\) равно 75.24.
    \item \(H_1: \mu \neq 75.24\) — математическое ожидание \(X_3\) отличается от 75.24 (двусторонняя альтернатива).
\end{itemize}

\subsubsection{Уровень значимости}
\(\alpha = 0.05\).

\subsection{Выбор критерия и обоснование}

\textbf{Выбранный критерий:} одновыборочный t-критерий Стьюдента для проверки гипотезы о математическом ожидании нормального распределения при неизвестной дисперсии (приложение A.5), реализованный функцией \texttt{scipy.stats.ttest\_1samp}.

Обоснование применимости:
\begin{itemize}
    \item Объём выборки \(n = 106\) достаточно велик; согласно центральной предельной теореме, распределение выборочного среднего близко к нормальному, что обеспечивает применимость t-критерия.
    \item Дисперсия генеральной совокупности неизвестна, поэтому используется её выборочная оценка.
\end{itemize}

\subsection{Вычисление статистики критерия}

Выборочные характеристики:
\[
\bar{x} = 74.725, \quad \hat{\sigma} = 10.553.
\]

Статистика критерия (приложение A.5):
\[
h = \sqrt{n} \, \frac{\bar{x} - \mu_0}{\hat{\sigma}}
= \sqrt{106} \cdot \frac{74.725 - 75.24}{10.553}
\approx -0.502.
\]

При справедливости \(H_0\) статистика \(h\) имеет распределение Стьюдента с \(\nu = n - 1 = 105\) степенями свободы.

\subsection{Принятие решения}

Согласно приложению A.3, используем p-value:
\[
p\text{-value} = 2 \cdot P(T_{105} \leq -0.502) \approx 0.616.
\]

\[
p\text{-value} = 0.616 > \alpha = 0.05.
\]

Следовательно, нет оснований отвергнуть \(H_0\).

\subsection{Вывод}

\textbf{Нет оснований отвергнуть нулевую гипотезу} \(H_0: \mu = 75.24\) на уровне значимости \(\alpha = 0.05\). Данные не противоречат предположению о том, что математическое ожидание \(X_3\) равно 75.24.

\subsection{Ошибка первого рода}

В данной задаче ошибка первого рода (приложение A.2) означает: \textbf{отвергнуть гипотезу \(H_0: \mu = 75.24\), когда на самом деле \(\mu = 75.24\)}. Вероятность такой ошибки \(\alpha = 0.05\).
